\documentclass[12pt]{article}

\usepackage[utf8]{inputenc}
\usepackage{amsmath, amssymb, amsthm}
\usepackage{geometry}
\geometry{a4paper, margin=2.5cm}

\title{Parametrizări și Primitivă în $\mathbb{C}$}
\author{}
\date{}

\theoremstyle{definition}
\newtheorem{definition}{Definiție}

\theoremstyle{plain}
\newtheorem{proposition}{Propoziție}

\begin{document}

\maketitle

\begin{definition}
(\textbf{Parametrizare în $\mathbb{C}$})  
O funcție 


\[
\gamma : [a,b] \to \mathbb{C}
\]


se numește \emph{parametrizare} dacă este continuă pe $[a,b]$ și derivabilă pe $(a,b)$.  
Punctele $\gamma(t)$ descriu traseul curbei în planul complex.
\end{definition}

\begin{definition}
(\textbf{Echivalența a două parametrizări})  
Două parametrizări 


\[
\gamma_1 : [a,b] \to \mathbb{C}, \qquad 
\gamma_2 : [c,d] \to \mathbb{C}
\]


sunt \emph{echivalente} dacă există o funcție bijectivă, continuă, strict monotonă și derivabilă


\[
\phi : [c,d] \to [a,b]
\]


astfel încât


\[
\gamma_2(t) = \gamma_1(\phi(t)) \quad \text{pentru toate } t \in [c,d].
\]


\end{definition}

\begin{definition}
(\textbf{Curbă închisă, netedă})  
O parametrizare $\gamma : [a,b] \to \mathbb{C}$ se numește \emph{curbă închisă} dacă


\[
\gamma(a) = \gamma(b).
\]


Ea este \emph{netedă} dacă $\gamma'$ este continuă pe $(a,b)$ și


\[
\gamma'(t) \neq 0 \quad \text{pentru toate } t \in (a,b).
\]


\end{definition}

\begin{definition}
(\textbf{Primitivă})  
Fie $f:\Omega \to \mathbb{C}$ o funcție continuă pe un domeniu $\Omega \subset \mathbb{C}$.  
O funcție $F:\Omega \to \mathbb{C}$ se numește \emph{primitivă} a lui $f$ dacă


\[
F'(z) = f(z) \quad \text{pentru toate } z \in \Omega.
\]


\end{definition}

\begin{proposition}
(\textbf{Leibniz–Newton pentru curbe în $\mathbb{C}$})  
Fie $f:\Omega \to \mathbb{C}$ o funcție care admite o primitivă $F$ în $\Omega$, și fie


\[
\gamma : [a,b] \to \Omega
\]


o parametrizare netedă.  
Atunci integrala lui $f$ de-a lungul curbei $\gamma$ este


\[
\int_{\gamma} f(z)\,dz = F(\gamma(b)) - F(\gamma(a)).
\]


În special, dacă $\gamma$ este o curbă închisă, atunci


\[
\int_{\gamma} f(z)\,dz = 0.
\]


\end{proposition}

\end{document}
